\section{Consequences for explanation, realism, and idealization}
\label{sec:scientific-methodology-consequences}

Observation--reflection equivalence changes the status of scientific explanation. An explanation is not merely an equation, derivation, mechanism, or model. It is a certified relation between finite observational access, extracted signatures, classifiers, stability conditions, and ledgered idealizations.

The physical reading follows the BEDC discipline for derived interfaces. Equality, function, number-like carriers, limits, real-like objects, and circle-like objects are not primitive simply because their names are familiar. They are derived interfaces licensed by certificates and ledgers. Likewise, a physical model is explanatory only when it shows how reality-facing signatures are preserved, compressed, stabilized, and made available for prediction.

\begin{definition}[Scientific explanation]
\label{def:scientific-explanation}
A scientific explanation of a phenomenon $P$ by a model $M$, relative to an observation bundle $\Pi$, is a certificate-mediated account showing that
$$
\operatorname{PhysSig}_{\Pi}(M)
\sim_{\Pi}
\operatorname{PhysSig}_{\Pi}(P),
$$
and that the predictive or explanatory operations used in the account descend through the approximation tower that produced $M$.
\end{definition}

\begin{definition}[Explanation certificate]
\label{def:explanation-certificate}
An explanation certificate is
$$
\operatorname{ExplCert}(M,P;\Pi)
=
(\operatorname{SourceSpec},
\operatorname{PhenomenonSpec},
\operatorname{PatternSpec},
\operatorname{ClassifierSpec},
\operatorname{StabilityCert},
\operatorname{DescentCert},
\operatorname{LedgerPolicy},
\operatorname{FailureSurface}).
$$
It records the observations from which the phenomenon is known, the phenomenon to be explained, the regularity or mechanism claimed, the comparison rule between model and reality, the stability basis for reuse, the descent proof for operations through the tower, the ledger of idealizations and residues, and the observations or refinements that would defeat the explanation.
\end{definition}

\begin{definition}[Explanatory scope and residue]
\label{def:explanatory-scope-residue}
The explanatory scope of $\operatorname{ExplCert}(M,P;\Pi)$ is the class of phenomena, histories, tower levels, and continuations over which the certificate remains valid. The explanatory residue $\operatorname{ExplRes}(P,M;\Pi)$ is the part of the phenomenon or system not captured by the certificate, including unmodelled variables, hidden tower levels, idealized-away structures, measurement conditions, and physical properties not queried by $\Pi$.
\end{definition}

\begin{theorem}[Explanation requires certificate]
\label{thm:explanation-requires-certificate}
A model explains a phenomenon only relative to an explanation certificate.
\end{theorem}

\begin{proof}
To explain a phenomenon, a model must be connected to observations of it, must state the pattern being explained, must compare model-generated and reality-generated signatures, must license reuse beyond the finite source trace, must show that explanatory operations descend through any compressions used, must record hidden source structure, and must specify the failure surface. These are exactly the fields of \autoref{def:explanation-certificate}.
\end{proof}

\begin{corollary}[Equations and derivations are not explanations by themselves]
\label{cor:equations-derivations-not-explanations}
An equation or derivation becomes explanatory only when embedded in a certificate specifying source, classifier, stability, descent, ledger, and failure surface.
\end{corollary}

\begin{proof}
An equation expresses a mathematical relation, and a derivation shows that a conclusion follows from assumptions internal to a model. Neither supplies the model-to-phenomenon certificate required by \autoref{thm:explanation-requires-certificate}.
\end{proof}

\begin{theorem}[Explanation cannot exceed scope]
\label{thm:explanation-cannot-exceed-scope}
A claim derived from $M$ is scientifically explained by $\operatorname{ExplCert}(M,P;\Pi)$ only when the claim lies within the certificate's explanatory scope.
\end{theorem}

\begin{proof}
The certificate licenses explanation only over its stated source traces, classifiers, tower levels, operations, stability domain, and continuations. Outside that scope at least one licensing field is missing or inactive.
\end{proof}

\begin{corollary}[Overextension is certificate failure]
\label{cor:overextension-certificate-failure}
A model is overextended when it is used beyond the scope of its explanation certificate.
\end{corollary}

\begin{proof}
By \autoref{thm:explanation-cannot-exceed-scope}, uses outside the certificate scope lack the needed source, classifier, stability, descent, or ledger field.
\end{proof}

\section{Certificate realism}
\label{sec:certificate-realism}

\begin{definition}[Certificate realism]
\label{def:certificate-realism}
Certificate realism is the position that realist commitment attaches to certified stable signature preservation, not to bare model syntax. A model $M$ is realist-worthy relative to $R$ and $\Pi$ when
$$
\operatorname{PhysSig}_{\Pi}(M)
\sim_{\Pi}
\operatorname{PhysSig}_{\Pi}(R)
$$
is supported by source, classifier, stability, descent, and ledger data.
\end{definition}

\begin{definition}[Signature realism and syntax realism]
\label{def:signature-realism-syntax-realism}
Signature realism is the special case of certificate realism in which the primary commitment is to stable invariant signatures generated by a system rather than to any single mathematical representation of those signatures. Syntax realism is the stronger claim that the mathematical syntax of a successful model directly describes the ontic structure of reality. The present framework treats syntax realism as under-certified unless the syntax is shown to preserve the relevant signatures and operations through the modelling tower.
\end{definition}

\begin{theorem}[Certificate realism lies between syntax realism and instrumentalism]
\label{thm:certificate-realism-between}
Certificate realism is weaker than syntax realism and stronger than instrumentalism.
\end{theorem}

\begin{proof}
It is weaker than syntax realism because it does not assert $M=R$; it asserts certified preservation of reality-facing signatures under specified bundles and tower levels. It is stronger than instrumentalism because it does not treat the model as merely useful. It claims that the model preserves stable signatures generated by the system, and those signatures can refute it.
\end{proof}

\begin{corollary}[Realism can be domain-local]
\label{cor:realism-domain-local}
A model may carry realist commitment within its certified signature domain while remaining open to later refinement.
\end{corollary}

\begin{proof}
Certificate realism is indexed by bundle, tower level, and continuation class. Later refinement may expose active residue without invalidating the certified commitment in the earlier domain.
\end{proof}

\begin{corollary}[Theory change is compatible with realism]
\label{cor:theory-change-compatible-realism}
If realist commitment attaches to certified signatures rather than full syntax, theory change need not destroy realism.
\end{corollary}

\begin{proof}
A later theory may replace earlier syntax while preserving the signatures certified by the earlier theory. The commitment transfers to the preserved signature content.
\end{proof}

\section{Effective theories}
\label{sec:effective-theories-methodology}

\begin{definition}[Effective theory]
\label{def:effective-theory-methodological}
An effective theory is a model $M_{\mathrm{eff}}$ whose explanation, prediction, and objectivity claims are valid only within a specified tower region
$$
\mathcal R_{\mathrm{eff}}
=
(\Pi,\mathcal C,L_i,\ldots,L_j),
$$
certified by source specification, coarse classifier, stability domain, descent domain, ledger policy, and failure surface.
\end{definition}

\begin{definition}[Effective exactness and domain leakage]
\label{def:effective-exactness-domain-leakage}
An effective theory is effectively exact over $\mathcal R_{\mathrm{eff}}$ when every admissible observation or operation inside that region preserves the relevant physical signatures. Domain leakage occurs when a claim is made at a level, bundle, or continuation class not included in $\mathcal R_{\mathrm{eff}}$ without an additional upgrade certificate.
\end{definition}

\begin{theorem}[Effective theories can be locally exact]
\label{thm:effective-theories-locally-exact}
An effective theory may be exact within its certified domain even if it is not globally final.
\end{theorem}

\begin{proof}
Effective exactness requires signature preservation for all admissible observations and operations inside the certified region. It does not require preservation under every possible refinement, scale, or future bundle.
\end{proof}

\begin{corollary}[Effective theories are not merely false]
\label{cor:effective-theories-not-merely-false}
An effective theory is not false merely because it fails outside its certified domain.
\end{corollary}

\begin{proof}
Failure outside the certified region shows that the certificate does not extend there. It does not defeat exactness inside the region licensed by \autoref{thm:effective-theories-locally-exact}.
\end{proof}

\begin{theorem}[Domain leakage requires upgrade certificate]
\label{thm:domain-leakage-requires-upgrade}
Using an effective theory outside its certified region is invalid unless an upgrade certificate extends the relevant source, classifier, stability, descent, and ledger fields.
\end{theorem}

\begin{proof}
A certificate licenses claims only inside its stated scope. Outside that scope the previous fields no longer apply. An upgrade certificate is precisely the extension of those fields to the new region.
\end{proof}

\begin{corollary}[Reduction is not automatic]
\label{cor:reduction-not-automatic-methodological}
A higher-level effective theory does not automatically reduce to a lower-level theory, and a lower-level theory does not automatically explain every higher-level structure.
\end{corollary}

\begin{proof}
Reduction across levels requires cross-level descent, signature preservation, and ledger compatibility. These are additional certificates, not consequences of level membership alone.
\end{proof}

\section{Idealization as licensed compression}
\label{sec:idealization-licensed-compression}

\begin{definition}[Idealization operator]
\label{def:idealization-operator}
An idealization operator
$$
I:L_i\rightsquigarrow L_j
$$
is a compression arrow replacing a lower-level source with a higher-level mathematical object, such as finite data by a continuous curve, a rough boundary by a smooth boundary, a molecular system by a continuum field, or a finite process by a limit object.
\end{definition}

\begin{definition}[Idealization certificate]
\label{def:idealization-certificate}
An idealization certificate for $I$ is
$$
\operatorname{IdealCert}(I)
=
(\operatorname{SourceSpec},
\operatorname{IdealPattern},
\operatorname{ClassifierSpec},
\operatorname{StabilityCert},
\operatorname{DescentCert},
\operatorname{ResidueLedger},
\operatorname{FailureSurface}).
$$
The residue ledger records what the idealization removes from the visible endpoint.
\end{definition}

\begin{definition}[Safe idealization]
\label{def:safe-idealization}
An idealization is safe relative to $\Pi$ when each removed residue is physically inactive under $\Pi$, or provably irrelevant by descent, or explicitly recorded as a failure surface.
\end{definition}

\begin{theorem}[Idealization requires certification]
\label{thm:idealization-requires-certification}
An idealized model is physically legitimate relative to $\Pi$ only if it carries an idealization certificate.
\end{theorem}

\begin{proof}
Idealization removes or compresses source structure. If the removed structure is inactive, that fact must be certified. If operations are applied to the ideal endpoint, they must descend through the compression. If removed structure may become active, it must be recorded as a failure surface. These are the fields of \autoref{def:idealization-certificate}.
\end{proof}

\begin{corollary}[Idealization hides rather than deletes]
\label{cor:idealization-hides-not-deletes}
When an idealization replaces a source by an ideal object, the hidden source structure remains part of the ledger.
\end{corollary}

\begin{proof}
\autoref{def:idealization-certificate} includes a residue ledger. The removed source structure is hidden at the endpoint but retained as provenance.
\end{proof}

\begin{corollary}[Smoothness is not free]
\label{cor:smoothness-not-free}
If smoothness is introduced by idealization, smoothness claims about the source require descent.
\end{corollary}

\begin{proof}
Smoothness may be introduced by interpolation, limiting, averaging, or completion. By \autoref{thm:idealization-requires-certification}, claims made using the smooth endpoint must descend through the idealization.
\end{proof}

\begin{theorem}[Idealization failure localizes residue]
\label{thm:idealization-failure-localizes-residue}
If an idealized explanation fails under a refined observation bundle, the failure identifies active residue hidden by the idealization.
\end{theorem}

\begin{proof}
A refined bundle that defeats the explanation exposes a distinction not preserved by the idealized endpoint. Since the endpoint was produced by hiding source structure, the exposed distinction belongs to the residue ledger.
\end{proof}

\section{Model pluralism and theory change}
\label{sec:model-pluralism-theory-change}

\begin{definition}[Certified model pluralism]
\label{def:certified-model-pluralism}
A family of models $\{M_\alpha\}_{\alpha\in A}$ exhibits certified model pluralism for $R$ relative to $\Pi$ when each $M_\alpha$ carries a physical-fit certificate for $R$ and all members preserve the same physical signature under $\Pi$:
$$
\operatorname{PhysSig}_{\Pi}(M_\alpha)
\sim_{\Pi}
\operatorname{PhysSig}_{\Pi}(M_\beta).
$$
\end{definition}

\begin{definition}[Translation certificate and representation equivalence]
\label{def:translation-certificate-representation-equivalence}
A translation certificate $\operatorname{TransCert}(M_1,M_2;\Pi)$ shows that the signatures, operations, ledgers, and failure surfaces of $M_1$ and $M_2$ correspond under $\Pi$. Two representations are representation-equivalent relative to $\Pi$ when such a certificate exists.
\end{definition}

\begin{theorem}[Model pluralism requires translation certificates]
\label{thm:model-pluralism-translation-certificates}
Multiple mathematical models can represent the same physical system relative to $\Pi$ when translation certificates show that they preserve the same physical signatures and compatible ledgers.
\end{theorem}

\begin{proof}
If each model preserves the physical signature of $R$ under $\Pi$, each is physically adequate there. If translation certificates also match their operations, ledgers, and failure surfaces, the models differ in representation rather than certified physical content.
\end{proof}

\begin{corollary}[Prediction agreement is not representation equivalence]
\label{cor:prediction-agreement-not-representation-equivalence}
Two models may agree on current predictions but fail to be representation-equivalent if their ledgers or failure surfaces differ.
\end{corollary}

\begin{proof}
Prediction agreement concerns present signatures. Representation equivalence requires translation of operations, ledgers, and failure surfaces.
\end{proof}

\begin{corollary}[Pragmatic model choice can remain realist]
\label{cor:pragmatic-choice-realist}
If two models are representation-equivalent under $\Pi$, choosing one for convenience does not undermine realism.
\end{corollary}

\begin{proof}
Realist commitment attaches to certified signatures, not syntactic form. If both models preserve the same certified signatures, convenience selects a representation, not the physical content.
\end{proof}

\begin{definition}[Theory transition]
\label{def:theory-transition}
A theory transition $T_1\leadsto T_2$ is a certified passage from one model or theory to another together with a map between their source specifications, classifiers, stability domains, ledgers, and failure surfaces.
\end{definition}

\begin{definition}[Conservative signature preservation and ledger refinement]
\label{def:conservative-signature-preservation-ledger-refinement}
A theory transition is conservative over $\Pi$ when every signature certified by $T_1$ under $\Pi$ is preserved by $T_2$. It is a ledger refinement when $T_2$ preserves old successful signatures while adding ledger information explaining where $T_1$ hid, approximated, or failed to descend.
\end{definition}

\begin{theorem}[Theory change can preserve realist content]
\label{thm:theory-change-preserves-realist-content}
If $T_1\leadsto T_2$ is conservative over $\Pi$, then the realist content certified by $T_1$ under $\Pi$ is preserved in $T_2$.
\end{theorem}

\begin{proof}
Under certificate realism, realist content is stable signature preservation. Conservative transition preserves those signatures, so the certified content remains present in the later theory.
\end{proof}

\begin{corollary}[Progress is ledger expansion]
\label{cor:progress-ledger-expansion}
Scientific progress can be understood as expansion and refinement of the ledger of preserved signatures, active residues, failure surfaces, and descent certificates.
\end{corollary}

\begin{proof}
A later theory improves an earlier one by preserving successful signatures, identifying failures, exposing residues, extending stability domains, and adding descent certificates. These are ledger refinements.
\end{proof}

\section{Failure surfaces and ideal objects}
\label{sec:failure-surfaces-ideal-objects}

\begin{definition}[Negative scientific knowledge]
\label{def:negative-scientific-knowledge}
Negative scientific knowledge is certified knowledge that a claim is not licensed under the current source, classifier, stability, descent, or ledger fields. It has the form
$$
\operatorname{CannotClaim}(\varphi;\mathcal C),
$$
where $\mathcal C$ records the missing or failed certificate field.
\end{definition}

\begin{definition}[Physical failure certificate]
\label{def:physical-failure-certificate}
A physical failure certificate for $\varphi$ is a typed witness that at least one required field of the relevant explanation, idealization, induction, or representation certificate fails, such as missing source, missing classifier, non-descent, absent ledger, stability failure, or unmatched signature.
\end{definition}

\begin{theorem}[Negative knowledge is scientific content]
\label{thm:negative-knowledge-scientific-content}
A certified refusal to assert $\varphi$ is part of the scientific content of a theory.
\end{theorem}

\begin{proof}
A theory is defined by its certification boundaries as well as by its positive claims. If $\varphi$ lacks source, classifier, stability, descent, or ledger support, asserting it would exceed the theory's scope. A certified refusal records that boundary.
\end{proof}

\begin{corollary}[Failure surfaces are explanatory]
\label{cor:failure-surfaces-explanatory}
A failure surface explains where and how a model may cease to apply.
\end{corollary}

\begin{proof}
It records the observations, refinements, or operations that would defeat the model's certificate, thereby identifying the limits of the explanation and the residues likely to become active.
\end{proof}

\begin{definition}[Certified ideal use]
\label{def:certified-ideal-use}
An ideal mathematical object has certified ideal use for a physical system when it carries a tower ledger, signature-preservation certificate, descent certificate, and failure surface.
\end{definition}

\begin{definition}[Ideal endpoint realism]
\label{def:ideal-endpoint-realism}
Ideal endpoint realism is the claim that the ideal object at the top of a modelling tower is itself the physical ontology. The present framework rejects this unless the endpoint is shown to preserve all relevant signatures and all operations used on it descend through the tower.
\end{definition}

\begin{theorem}[Ideal objects are admissible as certified endpoints]
\label{thm:ideal-objects-certified-endpoints}
An ideal mathematical object is physically admissible only when it is connected to physical source traces by a certified tower.
\end{theorem}

\begin{proof}
An ideal object has mathematical structure, but physical admissibility requires connection to reality-facing signatures. A certified tower supplies that connection by preserving signatures, recording hidden residue, and licensing operations by descent.
\end{proof}

\begin{corollary}[Usefulness does not imply ideal ontology]
\label{cor:usefulness-not-ideal-ontology}
An ideal mathematical object may be useful, predictive, and explanatory without being the ontic state itself.
\end{corollary}

\begin{proof}
Usefulness follows from certified ideal use. Ontic identity would require the stronger claim that the endpoint exhausts the ontic state, which \autoref{thm:ideal-objects-certified-endpoints} does not supply.
\end{proof}

\section{Scientific theory as certificate tower}
\label{sec:scientific-theory-certificate-tower}

\begin{definition}[Scientific theory]
\label{def:scientific-theory}
A scientific theory is a structured family
$$
\mathcal T_{\mathrm{sci}}
=
(\mathcal M,
\mathcal O,
\Pi,
\operatorname{FitCert},
\operatorname{LawCert},
\operatorname{ExplCert},
\operatorname{IdealCert},
\operatorname{Ledger},
\operatorname{FailureSurface}),
$$
where $\mathcal M$ is a family of models, $\mathcal O$ is a family of observed phenomena, $\Pi$ is a family of observation bundles, and the remaining fields certify fit, laws, explanations, idealizations, provenance, and failure conditions.
\end{definition}

\begin{definition}[Theory strength]
\label{def:theory-strength}
The strength of a scientific theory is determined by source breadth, classifier precision, predictive stability, descent depth, ledger completeness, failure-surface explicitness, and extent of inter-observer invariants. Theory strength is certificate strength.
\end{definition}

\begin{theorem}[Scientific theories are certificate-governed approximation towers]
\label{thm:scientific-theories-certificate-governed-towers}
Every physically serious scientific theory can be represented as a certificate-governed approximation tower.
\end{theorem}

\begin{proof}
A scientific theory begins from observations or experimental sources, extracts signatures, classifies them, packages them into models, stabilizes them into laws, uses them in explanations, and extends them into predictions. It also hides, idealizes, and compresses information, requiring ledgers and failure surfaces. This is exactly the structure of a certificate-governed approximation tower.
\end{proof}

\begin{corollary}[Theory comparison is certificate comparison]
\label{cor:theory-comparison-certificate-comparison}
Comparing two theories requires comparing their certificates, not merely their equations.
\end{corollary}

\begin{proof}
Theories may use different syntax while preserving the same signatures, or the same syntax while carrying different ledgers and failure surfaces. Source, classifier, stability, descent, ledger, and failure fields must therefore be compared.
\end{proof}

\begin{corollary}[A stronger theory may be less elegant]
\label{cor:stronger-theory-less-elegant}
A theory with more complete certificates may be scientifically stronger even if its mathematical form is less elegant.
\end{corollary}

\begin{proof}
Scientific strength is certificate strength. Elegance does not determine source coverage, classifier adequacy, stability, descent, ledger completeness, or failure-surface explicitness.
\end{proof}

\begin{remark}[Compressed form]
\label{rem:scientific-methodology-compressed-form}
The resulting image is:
$$
\text{scientific theory}
=
\text{models}
+
\text{signatures}
+
\text{classifiers}
+
\text{stability}
+
\text{descent}
+
\text{ledgers}
+
\text{failure surfaces}.
$$
Explanation is certified signature preservation; realism attaches to stable signatures; effective theories are tower-local; idealization is licensed compression; model pluralism requires translation certificates; and theory change is often ledger refinement rather than erasure.
\end{remark}

\concretizedIn{ch:concrete-instances-reality-constrained-explanation-namecert}{Reality-constrained explanations are realized as finite observation, model, phenomenon, signature-comparison, descent, residue, failure-surface, transport, route, provenance, and NameCert rows.}
